\chapter*{Заключение}                       % Заголовок
\addcontentsline{toc}{chapter}{Заключение}  % Добавляем его в оглавление

%% Согласно ГОСТ Р 7.0.11-2011:
%% 5.3.3 В заключении диссертации излагают итоги выполненного исследования, рекомендации, перспективы дальнейшей разработки темы.
%% 9.2.3 В заключении автореферата диссертации излагают итоги данного исследования, рекомендации и перспективы дальнейшей разработки темы.

Декодирование визуальной информации по сигналам мозга решено в три последовательных этапа, отвечающих возрастающей сложности учитываемой структуры данных: анализ временной динамики фМРТ, учет ее пространственно-временной структуры в условиях ограниченного объема выборки и мультимодальное объединение фМРТ и ЭЭГ. Анализ временной динамики показал, что корректное декодирование требует учета гемодинамической задержки BOLD-сигнала: ее оценка согласуется с известным физиологическим диапазоном и обеспечивает временное согласование стимула с фМРТ-откликом. На следующем этапе установлено, что пространственно-временные характеристики, выделяемые масками активности и проекцией ковариационных матриц на касательное пространство риманова многообразия, позволяют эффективно декодировать фМРТ в условиях ограниченного объема выборки, превосходя нейросетевые подходы. Завершающим этапом является предложенная мультимодальная архитектура, восстанавливающая визуальные стимулы по одновременным сигналам фМРТ--ЭЭГ: комбинированные эмбеддинги мозговой активности контрастивно выравниваются с CLIP-пространством, после чего генерация выполняется в две стадии~--- через априорную диффузионную модель и предобученную SDXL. Существенный методологический вклад~--- конвейер предобработки записей, который согласует модальности во времени, снижает размерность фМРТ и приводит ЭЭГ-сигналы к унифицированному набору каналов. Эксперименты подтверждают, что совместное использование модальностей превосходит одномодальные аналоги, а пространственно-временные характеристики данных фМРТ--ЭЭГ дополняют друг друга. Подключение диффузионной априорной модели дополнительно повышает семантическую согласованность реконструкций.

Кросс-доменный эксперимент дает ожидаемую картину: на стимулах, не представленных в обучающей выборке, качество падает~--- сказываются настройка энкодеров под конкретное распределение признаков и отсутствие встроенной доменной инвариантности. Этот результат очерчивает границы применимости метода и обозначает доменную адаптацию как направление дальнейшей работы. В целом предложенный подход закладывает прочную методологическую основу для последующих исследований в области мультимодального нейродекодирования и потенциально применим в задачах интерфейсов мозг--компьютер.

Основные результаты работы заключаются в следующем.

%% Согласно ГОСТ Р 7.0.11-2011:
%% 5.3.3 В заключении диссертации излагают итоги выполненного исследования, рекомендации, перспективы дальнейшей разработки темы.
%% 9.2.3 В заключении автореферата диссертации излагают итоги данного исследования, рекомендации и перспективы дальнейшей разработки темы.
\begin{enumerate}
    \item Задача декодирования визуальной информации по сигналам мозга формализована в единой постановке, объединяющей три уровня описания структуры данных: временную динамику, пространственно-временную структуру и мультимодальное объединение.
    \item Предложена линейная авторегрессионная модель прогнозирования фМРТ-отклика по визуальному стимулу; полученная оценка гемодинамического лага согласуется с известным физиологическим диапазоном.
    \item Разработан метод пространственно-временной классификации временных рядов фМРТ на основе масок активности и римановой геометрии ковариационных матриц, статистически значимо превосходящий нейросетевые аналоги в условиях ограниченной выборки.
    \item Предложена мультимодальная архитектура реконструкции визуальных стимулов по одновременным сигналам фМРТ и ЭЭГ; по метрике CLIP-Score она превосходит одномодальные аналоги, а априорная диффузионная модель дополнительно повышает семантическую согласованность реконструкций.
    \item Кросс-доменные эксперименты очертили границы применимости предложенных методов и обозначили доменную адаптацию как направление дальнейшей работы.
\end{enumerate}

